\section{热学}

\subsection{基础概念}
\begin{formulatab}{基础概念}
  体积（气体）&气体体积等于容器体积\\
  体积（液体）&液体体积由分子间距保持稳定\\
  压强定义&单位时间内对容器壁单位面积的平均冲击力\\
  温度定义&分子热运动的平均强度\\
\end{formulatab}

\textbf{微观解释}

体积、压强和温度都可以从分子热运动角度理解：气体分子无规则运动并不断碰撞容器壁，形成压强；温度反映分子热运动的剧烈程度。

\textbf{分子动理论}

物体由大量分子组成，分子在永不停息地做无规则运动，分子间同时存在引力和斥力。热现象的宏观规律由大量分子的统计行为决定。

\subsection{气体实验定律}
\begin{formulatab}{气体实验定律}
  玻意耳定律&$pV=C$\\
  查理定律（等压）&$p=CT$\\
  查理定律（等温）&$V=CT$\\
  理想气体状态方程&$pV=nRT$\\
\end{formulatab}

\textbf{应用说明}

气体实验定律描述了一定质量气体在不同约束条件下压强、体积和温度的关系。解题时通常先识别“等温、等压、等容”等过程，再选对应关系式。

\subsection{热力学定律}
\begin{formulatab}{热力学定律}
  热力学第一定律&$\Delta U=Q+W$\\
  等容过程&$W=0$\\
  等压过程&$W=p\Delta V$\\
  等温过程&$\Delta U=0$\\
  绝热过程&$Q=0$\\
\end{formulatab}

\textbf{第一定律}

第一定律给出内能变化与做功、热传递之间的关系。等容过程中系统不对外做功；等压过程中可用$W=p\Delta V$计算体积功；等温过程中理想气体内能不变；绝热过程中系统与外界无热交换。

\textbf{自由膨胀}

自由膨胀过程中气体失去外界压强约束，通常视为不做功。

\textbf{第二定律}

热量不能自发地从低温物体传向高温物体；不可能从单一热源吸收热量并全部转化为有用功。

\textbf{第三定律}

绝对零度无法达到。

\textbf{第零定律}

当系统之间停止热交换时，它们达到热平衡，温度相同。

\subsection{能量守恒与永动机}
\begin{formulatab}{能量守恒与永动机}
  能量守恒&能量既不会凭空产生，也不会凭空消失\\
  第一类永动机&本质上要求无中生有地产生能量\\
  第二类永动机&本质上要求热量100\%转化为有用功\\
\end{formulatab}

\textbf{能量守恒定律}

自然界中的能量可以在不同形式之间相互转化，但总量守恒。

\textbf{永动机不存在}

\begin{itemize}
  \item 第一类永动机：违反能量守恒。
  \item 第二类永动机：违反热力学第二定律。
\end{itemize}
